\section*{Chapitre \ref{ch:m3:when-markets-break-together} --- Quand les marchés se brisent ensemble}

\begin{solution}{exo:m3:when-markets-break-together:1}
$10\text{ millions}/0.125 = $ 80~millions de dollars.
\end{solution}

\begin{solution}{exo:m3:when-markets-break-together:2}
Dans une \omterm{def:m3:when-markets-break-together:flight}{fuite vers la qualité}, les investisseurs vendent les actifs risqués et achètent les plus sûrs, ce qui fait monter les cours des obligations d'État~; dans une ruée vers la liquidité, ils vendent tout ce qui peut
se vendre, actifs sûrs compris, pour lever des liquidités, et les cours des obligations d'État baissent avec le reste.
\end{solution}

\begin{solution}{exo:m3:when-markets-break-together:3}
Parce qu'il est détenu par les mêmes investisseurs à effet de levier qui ont besoin de liquidités~: dans une ruée vers la liquidité, les actifs sont vendus pour leur liquidité et non pour leurs perspectives, et
les corrélations estimées dans le calme se rompent.
\end{solution}

\begin{solution}{exo:m3:when-markets-break-together:4}
223.6~milliards de dollars contre 124.4~milliards sans impact sur les prix~: environ 99.3~milliards de dollars, soit 44~\%, proviennent des cours que font bouger les \omterm{def:m3:when-markets-break-together:fire}{ventes forcées}, et des pertes
qu'elles infligent.
\end{solution}

\begin{solution}{exo:m3:when-markets-break-together:5}
Environ $9 \times 0.64\% = 5.8\%$ de sa valeur en sept jours de négociation, pour l'actif le plus sûr du monde.
\end{solution}

\begin{solution}{exo:m3:when-markets-break-together:6}
La valeur en risque agrège les positions à l'aide d'une matrice de corrélation~; si les corrélations tendent vers un en situation de stress, les compensations estimées dans le calme disparaissent et les
pertes s'additionnent, si bien que l'estimation faite dans le calme sous-estime la perte précisément au moment où elle compte.
\end{solution}

\begin{solution}{exo:m3:when-markets-break-together:7}
Avec un effet de levier de 4, le choc ne crée aucune insuffisance de marge~: les fonds propres couvrent les marges plus élevées, rien n'est vendu, et la spirale ne démarre jamais. L'effet de levier décide
si le choc est absorbé ou amplifié.
\end{solution}

\begin{solution}{exo:m3:when-markets-break-together:8}
La diversification par classe d'actifs ne protège pas lorsque les détenteurs, le financement et les règles de marge sont communs~: dans un stress conjoint, les marges de la société augmentent
partout à la fois, et les actifs détenus par les mêmes investisseurs à effet de levier sont vendus ensemble.
\end{solution}

\begin{solution}{pb:m3:when-markets-break-together:1}
\textbf{1.} Environ 25.2~milliards de dollars. \textbf{2.} Environ 223.6~milliards de dollars. \textbf{3.} 8.88. \textbf{4.} 4.94, l'inverse du taux de marge moyen
après la hausse des marges (voir la \cref{prop:m3:when-markets-break-together:multiplier})~: chaque dollar vendu ne libère que le taux de marge exigé. \textbf{5.} Ils
baissent avec lui, à environ 93 et 80 le jour 20, car ils sont vendus pour faire face aux appels de marge sur l'ensemble du portefeuille. \textbf{6.} 8.73~: le résultat ne dépend
pas du pas de temps, car la réduction de levier est itérée jusqu'à un point fixe à chaque étape. \textbf{7.} Dans ce scénario, le choc est absorbé sans \omterm{def:m3:when-markets-break-together:fire}{vente forcée}~: la spirale
a besoin de marges qui augmentent. \textbf{8.} Les \omterm{def:m3:when-markets-break-together:fire}{ventes forcées} restent élevées, mais la corrélation des actifs non touchés avec l'actif choqué reste proche de zéro~: les détenteurs communs
sont ce qui relie les marchés entre eux. \textbf{9.} De 0.30 à 0.99. \textbf{10.} La \omterm{def:m3:when-markets-break-together:spirals}{spirale des marges} la déclenche ici, et l'impact sur les prix la double presque. \textbf{11.}
Les Treasuries sont vendus en même temps que les actions, les rendements grimpent brutalement, la profondeur de marché s'effondre et la base entre obligations au comptant et contrats à terme se disloque. \textbf{12.} Relèvement
des marges initiales d'environ 300~milliards de dollars sur le mois, les appels de marge de variation culminant à 140~milliards de dollars le 9 mars. \textbf{13.} Elle
contenait les swaps compensés~; la crise a transité par les dérivés non compensés et le financement à court terme, ainsi que les décotes sur collatéral de toute nature. \textbf{14.} Dans l'épisode
des gilts~: détenteurs à effet de levier, appels de marge croissant avec le mouvement, \omterm{def:m3:when-markets-break-together:fire}{ventes forcées} sur un marché en baisse, jusqu'à l'arrivée d'un acheteur en dernier ressort.
\textbf{15.} Les banques centrales ont acheté les actifs vendus et fourni des liquidités, ce qui rompt les deux spirales de l'extérieur. \textbf{16.} Elles doivent couvrir le risque
qu'elles perçoivent, mais des marges qui bondissent avec la volatilité sont procycliques~; des marges stables et prudentes, des coussins de liquidité constitués en période calme, ou des planchers
anti-procyclicité, réduisent la spirale. \textbf{17.} Il s'agit de la hausse des marges sur toutes ses positions à la fois, dans un stress calibré sur des épisodes comme celui de mars 2020, auquel s'ajoute
un financement qui risque de ne pas être reconduit, détenu en liquidités ou en actifs qui restent liquides dans une ruée vers la liquidité. \textbf{18.} Le levier et la concentration de ses positions sur ses marchés, les changements de marge
annoncés par les chambres de compensation, les spreads de financement et les décotes sur repo, la profondeur de marché et les corrélations entre des actifs habituellement indépendants.
\textbf{19.} \emph{Résultat nommé~:} un multiplicateur de \omterm{def:m3:when-markets-break-together:fire}{ventes forcées} de 8.88 dollars de \omterm{def:m3:when-markets-break-together:fire}{ventes forcées} par dollar d'insuffisance initiale de marge (8.73 avec huit pas
par jour), contre 4.94 sans impact sur les prix. \textbf{20.} Parce qu'ils partagent les détenteurs, le financement et les règles de marge, et que, dans une crise, ces liens propagent les ventes
d'un marché à l'ensemble des autres.
\end{solution}

\begin{solution}{iq:m3:when-markets-break-together:1}
Une \omterm{def:m3:when-markets-break-together:spirals}{spirale des pertes}~: les pertes forcent des ventes, les ventes font baisser les prix, et des prix plus bas causent des pertes. Une \omterm{def:m3:when-markets-break-together:spirals}{spirale des marges}~: la volatilité relève les marges, les marges forcent des ventes, les ventes
font monter la volatilité. Ensemble, elles transforment un choc en \omterm{def:m3:when-markets-break-together:fire}{ventes forcées} sur tout ce que détiennent les investisseurs à effet de levier.
\iqlookfor{les deux boucles de rétroaction et pourquoi elles traversent les marchés.}
\end{solution}

\begin{solution}{iq:m3:when-markets-break-together:2}
Dans la ruée vers la liquidité, les détenteurs ont vendu ce qu'ils pouvaient vendre pour lever des dollars~; les investisseurs non bancaires à effet de levier et les détenteurs étrangers ont vendu des Treasuries~; les dealers n'ont pas pu
absorber le flux, la profondeur s'est effondrée et les rendements ont grimpé.
\iqlookfor{la demande de liquidité, et non le risque de crédit.}
\end{solution}

\begin{solution}{iq:m3:when-markets-break-together:3}
À partir de fenêtres de stress passées plutôt que d'échantillons complets, en poussant les corrélations vers un pour les actifs détenus par les mêmes investisseurs à effet de levier, et par simulation
des \omterm{def:m3:when-markets-break-together:fire}{ventes forcées} sur les positions connues, vérifiée sur les épisodes de 2008, 2020 et 2022.
\iqlookfor{corrélations conditionnelles et raisonnement structurel.}
\end{solution}

\begin{solution}{iq:m3:when-markets-break-together:4}
Scénario~: hausse des marges égale à la pire hausse historique sur chaque plateforme à la fois, baisse conjointe des prix, lignes de financement non reconduites, \omterm{def:m3:blockchains-for-traders:stable}{décrochage} du \omterm{def:m3:blockchains-for-traders:stable}{stablecoin} (depeg) et retraits
des plateformes ralentis~; calculer le besoin de liquidités jour par jour au regard des liquidités et du collatéral liquide disponibles, délais de transfert entre plateformes compris~; fixer des coussins
et des limites pour que la société y survive sans \omterm{def:m3:when-markets-break-together:fire}{ventes forcées}.
\iqlookfor{marges conjointes, délais de transfert et échéancier de trésorerie jour par jour.}
\end{solution}

\begin{solution}{iq:m3:when-markets-break-together:5}
La base se creuse contre elle et les marges sur contrats à terme augmentent d'un coup~; avec un effet de levier de 50, un petit mouvement épuise ses fonds propres et elle doit vendre des Treasuries et racheter
les contrats à terme sur un marché en baisse, ce qui alimente la ruée vers la liquidité~; son prime broker peut relever les décotes ou retirer son financement.
\iqlookfor{marge, décote et dénouement forcé.}
\end{solution}

\begin{solution}{iq:m3:when-markets-break-together:6}
Des détenteurs avec positions et fonds propres, des règles de marge fondées sur la volatilité mesurée, une fonction d'impact sur les prix, un choc~; réduire le levier jusqu'à un point fixe à chaque pas, et exécuter
le calcul avec plusieurs tailles de pas, en vérifiant que les résultats convergent~; désactiver tour à tour chaque mécanisme et vérifier que les effets vont dans le bon sens.
\iqlookfor{points fixes, convergence selon le pas et ablations.}
\end{solution}
